\chapter{The theorem and its proof}
\label{ch:theorem}

Let $\mathcal Q_{222}\subset\mathbb R^8$ be the set of behaviours of
two separated quantum systems, with two binary measurements on each
system.  We use outcomes $\pm1$ and coordinates
\[
 q=(a_0,a_1,b_0,b_1,E_{00},E_{01},E_{10},E_{11}).
\]
The corresponding probabilities are
$p(\alpha,\beta\mid x,y)=
(1+\alpha a_x+\beta b_y+\alpha\beta E_{xy})/4$.
The quantum definition permits arbitrary finite local Hilbert-space
dimensions, mixed states, and binary positive-operator-valued measurements.

\begin{theorem}[Finite semidefinite representability]
\label{thm:main}
There are finite integers $N,m$ and real symmetric matrices
$A_0,A_1,\ldots,A_8,B_1,\ldots,B_m$ of order $N$ such that
\[
 \mathcal Q_{222}
 =\left\{q\in\mathbb R^8:
     \exists y\in\mathbb R^m,\quad
     A_0+\sum_{\nu=1}^{8}q_\nu A_\nu+
                 \sum_{j=1}^{m}y_jB_j\succeq0\right\}.
\]
\end{theorem}

The theorem is qualitative.  The argument below gives neither a value of
$N$ nor a bound of two on the orders of the positive-semidefinite blocks.
Its algebraic conclusion, which is the actual object proved, is stronger
than ordinary nonnegativity of every Bell slack.

\begin{theorem}[Separated slack certificate]
\label{thm:slack}
Let $K_c,K_t$ be separately labelled real closed extensions of
$\mathbb R$.  Let $f=f_0+\sum_{\nu=1}^{8} f_\nu q_\nu$ be an
affine function in $K_c[q]$ nonnegative on the transferred quantum set,
and let $q\in(\mathcal Q_{222})_{K_t}$.  Then
\[
 [f,q]:=f_0\otimes1+\sum_{\nu=1}^{8}f_\nu\otimes q_\nu
       \ \in\ \Sigma(K_c\otimes_{\mathbb R}K_t)^2.
\]
Each sum of squares is finite.  The fields may be labelled copies of
the same extension; no algebraic-disjointness hypothesis is imposed.
\end{theorem}

\section{A short proof, with every local input supplied below}

Write $P_{\mathcal Q}$ for the cone of valid affine functions.
The proof has four reductions.

\paragraph{1. A finite physical model.}
Chapter~\ref{ch:model} proves
\[
 \mathcal Q_{222}
   =\operatorname{conv}\Phi(S^1\times S^1\times S^3),
 \qquad
 q=\sum_{i=1}^{9}w_i\Phi(z_i),\quad
 w_i\geq0,\quad\sum_iw_i=1.
\]
Here $\Phi$ is the polynomial expectation map for two real qubits.
The common dilation, two-dimensional block decomposition, removal of
complex phases, and nine-term convexification are proved there.
Thus $\mathcal Q_{222}$ is compact, convex, and semialgebraic.
The fixed finite formula transfers to any real closed testing field.

\paragraph{2. A finite source reduction.}
Chapter~\ref{ch:geometry} proves that a valid source is an event sum
or a nonnegative combination of at most nine canonical sources.
Noncoercive sources have one of two explicit closed-face matrix lifts.
Every remaining extreme source has an actual finite stationary contact
and one of the following event supports, after a finite Bell relabelling:
\[
\begin{array}{ll}
\mathsf C_1:&\lambda=n=0,\quad d,m,C,A,\alpha>0,
                         \quad D,\beta\geq0,\quad D+\beta>0,\\
\mathsf C_2:&\lambda=\beta=0,\quad d,m,C,A,D,\alpha>0,
                         \quad n\geq0.
\end{array}
\]
The geometry uses the nonnegative outer coefficients of the contact
determinant, an eight-event compatibility calculation, and the six
linear contact equations.  It does not infer a tensor certificate from
a determinant certificate.  Chapter~\ref{ch:entry} proves the strict
native signs and identifies the original stationary coefficients.

\paragraph{3. The complete canonical polynomial.}
For $D_u=1+u^2$, $D_v=1+v^2$ and $\psi=(W,X,Y,Z)$, let
\[
\begin{split}
F={}&4\lambda D_u(W+vX)^2+4nD_uD_vX^2
       +4m(vW-X+uvY-uZ)^2\\
 &+4CD_v(uW-Y)^2+4AD_uD_vY^2
       +(4D+2d)(uW+uvX-Y-vZ)^2\\
 &+(4\alpha+2d)D_u(vY-Z)^2
       +(4\beta+2d)D_v(X+uZ)^2-2dD_uD_vZ^2.
\end{split}                                                     \tag{M1}
\]
The canonical certificate theorem in Chapter~\ref{ch:cover} proves
$F\in\Sigma(K_c\otimes K_t)^2$ for both supports, every independent
testing pair, and every state.  Its proof is the finite source cover
in that chapter together with the full identities in the certificate
chapters.  Those identities include the scalar terms, both independent
state normals, and the remaining state quadratic form.  They retain
literal zero event weights, arbitrary relative infinitesimal rates,
all state ratios, and the homogeneous angle endpoints.

There is one algebraic mechanism behind these certificates.  If
$G$ has coefficients in the tensor product of the two bounded rings
and has a positive-definite real residue, then $z^TGz$ is tensor SOS.
Chapter~\ref{ch:tensor} proves this by a finite elementary-square
identity.  The roots $z$ may be unbounded.  The common contact
division and circuit lemmas in Chapter~\ref{ch:mechanisms} produce
the exact bounded coefficients; the certificate chapters prove the
strict residue and cover the testing complements.  No mixed inverse
is introduced in this step.

\paragraph{4. Restore the slack and assemble.}
The homogeneous version of (M1) equals
\[
 (1\otimes\Delta_A\Delta_B\|\psi\|^2)\,[f,\Phi(z)],
 \qquad \Delta_A,\Delta_B>0.
\]
Every cancelled factor is a nonzero pure-source or pure-testing
factor.  The closed-face lifts give separated trace-pairing squares;
events give immediate elementary squares.  For a source decomposition
$f=\sum_j a_j f_j$ and a testing mixture $q=\sum_iw_iq_i$, write
$[f_j,q_i]=\sum_k h_{ijk}^2$.  Then, exactly,
\[
 [f,q]=\sum_{i,j,k}
       \bigl[(\sqrt{a_j}\otimes\sqrt{w_i})h_{ijk}\bigr]^2.
                                                               \tag{M2}
\]
The constant coefficient is preserved by $\sum_iw_i=1$; zero
weights cause no division.  This proves Theorem~\ref{thm:slack}.

Finally apply Scheiderer's tensor-evaluation criterion
\cite[Corollary 3.19]{scheiderer} to two labelled copies of each
arbitrary real closed field $K$.  Its hypotheses are exactly
closedness, convexity, semialgebraicity, and the separated SOS
assertion just established.  It yields a finite semidefinite lift,
which is Theorem~\ref{thm:main}.\hfill$\square$

\section{Reading the certificate chapters}

The short argument above is the proof's logical spine.  The complete
proof consists of that argument and the explicit certificates that
follow.  The latter are essential: replacing them by a list of names
would leave the canonical certificate theorem unproved.

Three conventions keep the detailed calculations unambiguous.
First, coefficients and actual contact parameters always belong to
$K_c$; independent testing angles and state amplitudes belong to
$K_t$.  Second, $\mathrm I$, $\mathrm U$, and $\mathrm O$ mean,
respectively, positive infinitesimal, positive bounded with positive
real residue, and positive unbounded.  These letters refer to the
field specified in each statement.  A letter H in an inherited chart
name means type $\mathrm O$, not a literal infinite coordinate.
Third, a displayed equation tag such as RC or SC is local to its
certificate section.  Section and page references identify the
intended formula when several families use the same short tag.

The refinement removes superseded proof routes and consolidates
shared algebra while preserving the formulas on which the remaining
routes depend.  It is a reduction of the recovered argument, not a
claim that no shorter proof can exist.  A separate verification record
in the source package identifies the new exact algebra checks and
the substantive proof replacements.
